\chapter{Introduction}
\label{ch:intro}

\section{The clinical problem}

Oesophageal adenocarcinoma has a five-year survival below 20\% in most Western series, and almost all cases arise from Barrett's oesophagus, a metaplastic change in the lower oesophagus that is surveilled by periodic endoscopy and biopsy \citep{fitzgerald2014bsg}. Surveillance generates two things in large quantities: H\&E-stained biopsy slides and free-text pathology reports grading them on the ladder non-dysplastic Barrett's (NDBE), indefinite for dysplasia (IND), low-grade dysplasia (LGD), high-grade dysplasia (HGD) and cancer. The annual progression rate from NDBE is well under 1\%, so the central clinical question is which of the many patients under surveillance will progress, and the central data problem is that the outcome is rare, the follow-up is long, and the labels live in prose.

Two families of computational work address this. Genomic studies showed that copy-number changes measured by shallow whole-genome sequencing (sWGS) of biopsies predict progression years ahead of the histological diagnosis \citep{killcoyne2020}; commercial tissue-systems assays such as TissueCypher stratify risk from multiplexed immunofluorescence \citep{critchleythorne2016}. Deep learning on whole-slide images has, separately, reached clinical-grade performance on diagnostic tasks when trained on tens of thousands of slides labelled only from their reports \citep{campanella2019}, and pathology foundation models now provide strong tile-level features without task-specific training \citep{chen2024uni, xu2024gigapath, vorontsov2024virchow}.

The obvious next step, and the one this thesis started from, is to fuse the two: use H\&E as the anchor modality that every patient has, add genomic or clinical data where it exists, and expect the combination to predict progression or survival better than either alone. Pan-cancer studies on TCGA report exactly this \citep{chen2022porpoise}.

\section{What this thesis found instead}

The fusion gain did not replicate in oesophageal cancer, and working out why turned into the thesis. Chapter~\ref{ch:fusion} reports the replication attempts: four cohorts, seven fusion architectures, four tile encoders and a published multimodal baseline. Every fusion arm's advantage over the best single modality has a confidence interval that includes zero, once three checks are applied that the published literature mostly omits. The first is selection: when the best of several fusion arms is reported, its advantage must be tested against the distribution of the \emph{best} arm's advantage under the null, not the distribution of one arm's. The single result that had survived a pre-registered, Holm-corrected confirmatory test did not survive this check. The second is patient overlap: two cohorts treated as independent turned out to share 36\% of one cohort's patients, and a transfer result that had looked like modest generalisation collapsed to chance when they were removed. The third is prevalence: pooling cohorts with different mutation rates lets a model score well by learning which cohort a slide came from, and a predictor given only cohort identity matched the pooled probes.

Alongside these negatives, a simulation-based power map (Section~\ref{sec:power}) shows that none of the available cohorts can detect a fusion gain smaller than 0.075 to 0.10 in AUC or concordance at 80\% power, whereas the gains actually observed are +0.01 to +0.04. The honest statement is therefore not that fusion does not help in oesophageal cancer but that no cohort of this size can show that it does, and that several published-style analyses would have declared a win here for reasons that have nothing to do with fusion.

The second half of the thesis is about supervision. Report-derived labels are the only way to reach the slide volumes at which histology models work, and Chapter~\ref{ch:labels} builds and validates a labelling pipeline that runs entirely inside the hospital network: a jury of eight open-weight language models, with an explicit unsure class, human adjudication of the hardest cases, and validation against human registry grades in cancer types the jury had never seen. Re-labelling at the level of individual specimen sections, prompted by the observation that a single report often grades several biopsies differently, quantified for the first time how much label noise the usual report-level shortcut introduces (32\% of slides), and then showed that the shortcut is nevertheless no worse for training. Chapter~\ref{ch:vlm} takes the same report-slide pairs and trains a contrastive vision-language model that grades slides without any labels, and reports where that model does and does not transfer.

\section{Contributions}

\begin{enumerate}
\item A pre-registered, four-cohort, multi-encoder replication of histology-plus-genomics fusion in oesophageal cancer, with a uniformly null result once selection, overlap and prevalence are controlled (Chapter~\ref{ch:fusion}).
\item A power map for fusion gains in cohorts of the sizes available in this disease, and the argument that failure to replicate is mostly failure to power (Section~\ref{sec:power}).
\item A six-rule protocol for claiming a fusion gain, each rule motivated by a specific failure in this thesis (Chapter~\ref{ch:discussion}, Appendix~\ref{app:p2}).
\item A locally-run language-model jury for pathology reports, with robustness, uncertainty and external validation results, and a quantification of section-level versus report-level label noise (Chapter~\ref{ch:labels}, Appendix~\ref{app:p1}).
\item A contrastive slide-report model with zero-shot grading at AUC 0.889 and a negative transfer result (Chapter~\ref{ch:vlm}).
\item A documented research process: a public execution plan with 46 numbered experiments, dated pre-registrations and deviations, and four waves of adversarial review that changed three conclusions (Appendices~\ref{app:plan} and \ref{app:review}).
\end{enumerate}

\section{How to read this thesis}

Chapter~\ref{ch:cohorts} describes the cohorts and the audit of their independence. Chapter~\ref{ch:methods} fixes the evaluation machinery shared by every experiment: patient-disjoint folds, paired bootstraps, permutation controls, the multiplicity plan and the pre-registration contract. Chapters~\ref{ch:fusion} to \ref{ch:vlm} report results. Chapter~\ref{ch:discussion} draws them together. The appendices hold the execution plan summary, two paper drafts, the review process and a register of every numbered claim with its evidence file. All numbers in the text are transcribed from JSON result files committed to the public repository (\url{https://github.com/rzuberi/thesis}); the register in Appendix~\ref{app:claims} lists which file backs which claim. No patient-level data or report text appears in the repository or in this document.
